\documentclass[12pt,a4paper]{article}

\usepackage{fontawesome5}
\providecommand{\faPencilAlt}{\faIcon{pencil-alt}}

\usepackage{styles/couleurs}
\usepackage{styles/mise_en_page}
\usepackage{styles/boites}
\usepackage{styles/maths}
\usepackage{styles/figures}

\usepackage{enumitem}
\usepackage{array}
\usepackage{tabularx}
\usepackage{amsmath}
\usepackage{amssymb}
\usepackage{graphicx}
\usepackage{tikz}
\usepackage{xparse}

\usetikzlibrary{calc,arrows.meta}

\setlist[itemize]{leftmargin=1.4em,itemsep=0.15em,topsep=0.2em}
\setlist[enumerate]{leftmargin=*,itemsep=0.22em,topsep=0.25em}

\newcommand{\lignepointillee}{%
  \noindent\dotfill\par
  \vspace{0.45em}
}

\newcommand{\lignesreponse}[1]{%
  \begin{tcolorbox}[
    enhanced,
    colback=white,
    colframe=NeutreC!40,
    boxrule=0.5pt,
    arc=1.5mm,
    left=3mm,
    right=3mm,
    top=2mm,
    bottom=1mm,
    before skip=2pt,
    after skip=2pt
  ]
    \foreach \i in {1,...,#1}{\lignepointillee}
  \end{tcolorbox}
}

\newcommand{\imageactivite}[3]{%
  \IfFileExists{#3}{%
    \includegraphics[width=#1,height=#2,keepaspectratio]{#3}%
  }{%
    \fbox{%
      \parbox[c][#2][c]{#1}{%
        \centering\sffamily\small
        Image à intégrer\par
        \vspace{0.3em}
        \texttt{\detokenize{#3}}
      }%
    }%
  }%
}

\newcommand{\titreseq}[1]{%
\begin{center}
  {\sffamily\LARGE\bfseries\color{TitreC}
  Séquence 5 -- Statistiques et proportionnalité}\\[0.2em]
  {\sffamily\large\bfseries\color{SavoirC}#1}
\end{center}
}

\begin{document}

\titreseq{Automatismes -- Séries courtes}

\remarque{Chaque série est prévue pour environ cinq minutes.
On peut projeter une seule série au début d'une séance.}

\section*{Série 1 -- Grandeurs et unités}

Pour chaque situation, nomme les deux grandeurs et leurs unités.

\begin{enumerate}
  \item \(4\) cahiers coûtent \(12\) €.
  \item Une voiture parcourt \(150\) km en \(2\) h.
  \item Une recette utilise \(300\) g de farine pour \(4\) personnes.
\end{enumerate}

\section*{Série 2 -- Proportionnel ou non ?}

\begin{enumerate}
  \item Nombre de stylos identiques et prix total.
  \item Âge d'un enfant et âge de son père.
  \item Côté d'un carré et périmètre.
  \item Côté d'un carré et aire.
\end{enumerate}

\section*{Série 3 -- Compléter une colonne}

\[
\begin{array}{c|ccc}
\text{Quantité} & 3 & 5 & 8\\
\hline
\text{Prix en €} & 7{,}5 & & 20
\end{array}
\]

\section*{Série 4 -- Retour à l'unité}

\[
6\text{ kg}\longleftrightarrow15\text{ €}
\]

Calcule le prix de \(1\) kg, puis de \(4\) kg.

\section*{Série 5 -- Coefficient}

Dans chaque tableau, trouve le coefficient de la première ligne vers la seconde.

\[
\begin{array}{c|ccc}
 & 2 & 5 & 9\\
\hline
 & 7 & 17{,}5 & 31{,}5
\end{array}
\qquad
\begin{array}{c|ccc}
 & 4 & 7 & 10\\
\hline
 & 1 & 1{,}75 & 2{,}5
\end{array}
\]

\section*{Série 6 -- Effectifs et fréquences}

Dans une classe de \(25\) élèves, \(10\) viennent à pied.

Calcule la fréquence puis le pourcentage.

\section*{Série 7 -- Pourcentages et angles}

Complète mentalement.

\[
50\%\longleftrightarrow\ldots^\circ
\qquad
25\%\longleftrightarrow\ldots^\circ
\]

\[
10\%\longleftrightarrow\ldots^\circ
\qquad
20\%\longleftrightarrow\ldots^\circ
\]

\section*{Série 8 -- Contrôle d'un diagramme}

Un diagramme annonce \(35\%\), \(28\%\), \(22\%\) et \(20\%\).

Peut-il représenter une répartition complète ? Justifie.

\end{document}
